\section{The Sacraments of the Departing}

The craft of dying is not left to the devotion of the dying man and his friends. Christ gave his Church sacraments for this passage, and she gives them as long as there is breath to receive them. ``Penance, the Anointing of the Sick and the Eucharist as viaticum constitute at the end of Christian life `the sacraments that prepare for our heavenly homeland' or the sacraments that complete the earthly pilgrimage.''\endnote{CCC 1525.}

The first is confession. The \emph{Craft of Dying} told the sick man to make his peace with God ``before all other things,'' and to seek his soul's health ``by very contrition and confession.''\endnote{\emph{Craft of Dying}, ch.~5, Comper, 33--34.} At the end the Church's power to bind and loose is offered once more, and the soul that has strayed far may return in an hour. Those who have confessed each month ``as if it were the last of your life,'' as Leonard of Port Maurice counsels, come to this confession as to one long rehearsed.

The second is the anointing that the Apostle James commands: ``Is any man sick among you? Let him bring in the priests of the church and let them pray over him, anointing him with oil in the name of the Lord. And the prayer of faith shall save the sick man. And the Lord shall raise him up: and if he be in sins, they shall be forgiven him.''\endnote{James 5:14--15, Douay--Rheims (Challoner).} The Council of Trent defines that this anointing is truly and properly a sacrament, instituted by Christ and promulgated by the Apostle.\endnote{Council of Trent, session XIV (1551), canons on extreme unction, canon 1, DS 1716.} In the Roman Ritual the priest anointed the senses through which the man had sinned, and at each he invoked the holy unction and the most tender mercy of God; at the lips he prayed, ``may the Lord pardon thee whatsoever sins thou hast committed by taste and speech.''\endnote{\emph{Rituale Romanum} (1872), the form of extreme unction at the mouth, \latin{indulgeat tibi Dominus quidquid per gustum et locutionem deliquisti}; English from \emph{A Manual of Prayers} (1889), 488.} The present law of the Latin Church allows the sacrament to be given to any of the faithful ``who, having reached the use of reason, begins to be in danger due to sickness or old age,'' and to be repeated if the danger returns or grows.\endnote{\emph{Code of Canon Law} (1983), can.~1004, official English (Vatican web text).} The Church does not wait for the last breath to strengthen the sick against the last struggle.

The third is the Body of the Lord given as food for the journey. When the priest brings Holy Communion to one who is dying, he says over him words the Church has spoken at countless bedsides: \latin{Accipe, frater (vel soror) Viaticum Corporis Domini nostri Jesu Christi, qui te custodiat ab hoste maligno, et perducat in vitam \ae ternam. Amen.} ``Receive, brother, the Viaticum of the Body of our Lord Jesus Christ; may He preserve thee from the wicked enemy, and bring thee unto life everlasting. Amen.''\endnote{\emph{Rituale Romanum} (1872), De communione infirmorum; English from \emph{A Manual of Prayers} (1889), The Holy Viaticum, 478.} The Catechism calls this the last sacrament of all: ``As the sacrament of Christ's Passover the Eucharist should always be the last sacrament of the earthly journey, the `viaticum' for `passing over' to eternal life.''\endnote{CCC 1517.} Christ himself promised it: ``He who eats my flesh and drinks my blood has eternal life, and I will raise him up at the last day.''\endnote{John 6:54, as quoted in CCC 1524.}

Because the Viaticum matters so much, the Church's law is unusually urgent about it. ``Holy Viaticum for the sick is not to be delayed too long; those who have the care of souls are to be zealous and vigilant that the sick are nourished by Viaticum while fully conscious.''\endnote{\emph{Code of Canon Law} (1983), can.~922; can.~921 directs that the faithful in danger of death from any cause be nourished by Holy Communion in the form of Viaticum.} The canon asks something of the dying and of those who love them. It asks the sick man not to put off the priest for fear of frightening himself, and it asks his family not to put him off for fear of frightening the sick man. Love at a deathbed sends for the priest in time.

A Christian who has received these three sacraments at the end has been given everything the Church has to give. He has been forgiven, strengthened, and fed. He goes out of the world as he entered the Church, carried by grace, and the Church's last service is to accompany him with her prayer to the door.
